\section{Giriş}
\label{sec:giris}

Otonom ay inişi, güç inişi (powered descent) evresinde gerçek zamanlı
olarak kraterli, kayalık ve eğimli bir arazi üzerinde hassas
konumlandırma ve güvenli temas gerektiren, doğrudan operasyonel sonuçları
olan bir kontrol problemidir. Son yıllarda GPU paralel fizik motorları
(NVIDIA PhysX) ve buna dayalı simülasyon çerçeveleri, binlerce ortamı eş
zamanlı çalıştırarak derin pekiştirmeli öğrenmeyle (RL) karmaşık temas
zengini görevlerin pratik sürelerde öğrenilmesini mümkün kılmıştır
\citep{mittal2023orbit,nvidia2025isaaclab,rudin2022walking}. Bu
çalışmanın temelini oluşturan proje de bu sınıfa girer: dört bacaklı, iki
eksenli gimbal itkili bir roketin, Soft Actor-Critic (SAC) algoritmasıyla,
gerçek dijital yükseklik modeli (DEM) tabanlı bir makro yüzey üzerine
bindirilmiş kraterler, kayalar ve santimetre ölçekli regolit dokusuyla
zenginleştirilmiş prosedürel bir ay zemininde inişi öğrendiği bir Isaac
Lab doğrudan iş akışı (direct-workflow) görevidir.

Bu tarz prosedürel çok çözünürlüklü arazilerde, hesaplama ve bellek
bütçesini sınırlı tutmak için tipik olarak iki katman kullanılır: geniş
alanı kapsayan düşük çözünürlüklü bir küresel (global) yüzey ve ajanın ya
da hedefin yakınında yalnızca gerekli bölgede oluşturulan yüksek
çözünürlüklü bir yerel (local) yama. Bu iki katmanın dikişsiz birleşmesi,
bilgisayar grafikleri literatüründe yirmi beş yılı aşkın süredir
çalışılan klasik bir problemdir (geometry clipmaps için
\citep{losasso2004geometry}; ROAM için \citep{duchaineau1997roaming};
geomipmapping için \citep{deboer2000geomipmapping}). Ancak bu literatür
neredeyse tamamen görsel işleme (rendering) kalitesine odaklanır;
birleşme sınırındaki bir süreksizliğin, ayağın zeminle temasını, yükseklik
sensörü okumalarını veya bir RL ajanının proprioseptif gözleminin
gürültü karakterini nasıl etkilediği nicel olarak nadiren ele alınır.

Bu problem, incelenen sistemde iki farklı ve birbiriyle tutarsız çözümle
karşımıza çıkmaktadır: eğitim ortamı katmanlar arasında geçişi kompakt
destekli bir kübik smoothstep fonksiyonuyla sağlarken, önceden kullanılan
mesh tabanlı bir uygulama sonsuz destekli bir Gauss çanı kullanmaktadır.
Bu iki fonksiyon, matematiksel olarak temelden farklı bir özelliğe
sahiptir: smoothstep sonlu bir yarıçapta tam olarak sıfıra iner (kompakt
destek), oysa standart Gauss çanı hiçbir sonlu yarıçapta tam sıfıra
inmez (sonsuz destek; \citep{hardy1971multiquadric}) ve bir mühendislik
uygulamasında kullanılabilmesi için mutlaka bir noktada kesilmesi
(truncation) gerekir; bu kesme işlemi, kesim noktasında yeni ve
ölçülebilir bir dikiş (seam) yaratma riski taşır. Bu iki yöntemi adil
biçimde karşılaştırmak, ikisini de aynı geçiş yarıçapında sıfıra
zorlamayı (hard truncation) ve aynı geçiş bütçesi altında ölçmeyi
gerektirir; bu karşılaştırma biçimi, yapılan literatür taramasında
terrain-randomization çalışmalarının hiçbirinde doğrudan ele alınmamıştır
(bkz. \cref{sec:ilgili}).

Bu gözlemlerden hareketle çalışma üç araştırma sorusuna yanıt aramaktadır:

\begin{itemize}
\item \textbf{AS1:} Kompakt destekli (smoothstep) ve sonsuz destekli
(Gauss) karıştırma fonksiyonları, aynı geçiş yarıçapı bütçesi altında
adil biçimde karşılaştırıldığında, sınır sürekliliği (C0/C1) açısından
nicel olarak ne kadar farklıdır ve bu fark hesaplama/bellek maliyetiyle
nasıl bir denge (trade-off) oluşturur?
\item \textbf{AS2:} İniş ayağı temas bölgesinde Gauss, ötesinde
smoothstep kullanan bir hibrit karıştırma yöntemi, her iki yöntemin
süreklilik avantajlarını makul bir hesaplama maliyeti artışıyla
birleştirebilir mi?
\item \textbf{AS3:} Bu üç karıştırma yönteminin sürekliliğindeki farklar,
gerçek bir Isaac Sim/Isaac Lab SAC eğitiminde ölçülebilir bir öğrenme
performansı farkına yansır mı?
\end{itemize}

Bu doğrultuda çalışmanın amacı üç yönlüdür: (i) mevcut smoothstep
yöntemini ve önceden kullanılan Gauss yöntemini aynı mimaride, aynı
geçiş bütçesi altında nicel olarak karşılaştırmak; (ii) iki yöntemin
avantajlarını birleştiren bir hibrit yöntem önerip doğrulamak; (iii) bu
üç yöntemi yalnızca analitik/CPU düzeyinde değil, gerçek bir Isaac
Sim/Isaac Lab GPU paralel fizik simülasyonu içinde, gerçek bir SAC
eğitimiyle karşılaştırarak, karıştırma fonksiyonu seçiminin öğrenme
performansına yansıyıp yansımadığını göstermektir. Bu son nokta,
çalışmanın özgün değerinin merkezindedir: karşılaştırma soyut bir sayısal
deney olarak değil, gerçek ay zemini üzerinde, gerçek bir Isaac Sim
eğitim döngüsü içinde gerçekleştirilmiştir.

Bu çalışmanın temel katkıları şunlardır:

\begin{itemize}
\item Var olan smoothstep ve önceden kullanılan Gauss karıştırma
yöntemleri, sonlu geçiş yarıçapında kesilerek (hard truncation) aynı
geçiş bütçesi altında ilk kez adil biçimde karşılaştırılmış; kesilmemiş
Gauss'un görünürde iyi durduğu ama gerçekte sonlu yarıçapta kesildiğinde
ortaya çıkan gizli bir dikiş kusuru taşıdığı gösterilmiştir.
\item Temas bölgesinde Gauss, ötesinde smoothstep ağırlıklı bir switch
kullanan yeni bir hibrit karıştırma yöntemi önerilmiş ve bu yöntemin
makine hassasiyetinde süreklilik sağladığı, kendi iç geçişinde yeni bir
dikiş yaratmadığı gösterilmiştir.
\item Otuz tohum ve dört yöntem varyantı (smoothstep, üretim öncesi
Gauss, kalibre edilmiş Gauss, hibrit) üzerinde istatistiksel olarak
doğrulanmış (Wilcoxon işaretli sıra testi) bir süreklilik, maliyet ve
VRAM karşılaştırması yapılmış; dört ayrı ablasyon analiziyle (temas
yarıçapı, switch genişliği, geçiş yarıçapı, grid çözünürlüğü)
bulguların sağlamlığı sınanmıştır.
\item Üç yöntem, gerçek bir Isaac Lab/Isaac Sim PhysX ortamında, GPU
paralel SAC eğitimiyle karşılaştırılmıştır; bu, terrain-randomization
literatüründeki çalışmaların aksine karıştırma fonksiyonu seçimini
doğrudan bir eğitim tasarımı değişkeni olarak ele alan ilk denemedir.
\item Bu gerçek eğitim karşılaştırması yöntem başına tek tohumla değil
üç bağımsız tohumla ($n=3$) yapılmış; ilk tek-tohumlu ölçümde görünen
büyük performans farkının çoklu tohumla istatistiksel olarak
doğrulanamadığı gösterilerek, tek-koşu RL karşılaştırmalarının
güvenilirliği konusunda somut, kendi verisine dayanan bir uyarı örneği
sunulmuştur.
\item Üç yöntem de üretim koduna, mevcut davranışı bozmayan (geriye
dönük uyumlu) ve bir ortam değişkeniyle
(ISAACLAB\_TERRAIN\_BLEND\_METHOD) seçilebilir şekilde entegre
edilmiştir.
\end{itemize}

Bu makalenin geri kalanı şu şekilde yapılandırılmıştır: \cref{sec:ilgili},
çalışmanın konumlandığı literatürü beş tematik küme altında
incelemektedir. \cref{sec:yontem}, deney ortamını ve tüm koşuların
paylaştığı simülasyon kurulumunu, karıştırma fonksiyonlarının matematiksel
tanımını, analitik değerlendirme protokolünü ve gerçek Isaac Sim eğitim
protokolünü ayrıntılandırmaktadır. \cref{sec:bulgular}, süreklilik,
maliyet, VRAM, ablasyon ve gerçek eğitim sonuçlarını sunmaktadır.
\cref{sec:tartisma}, bulguları literatürle ilişkilendirerek tartışmakta ve
sınırlamaları ele almaktadır. \cref{sec:sonuc}, çalışmayı özetlemekte ve
somut gelecek çalışma önerileriyle sonlanmaktadır.
